% ==============================================================
% Diskussionspapier
%
% Zeit, Entropie und informationelle Transitionen
%
% Vorbereitung für das Modul:
% Die Projektive Chronogenese (DPC)
%
% Version 2.0 (Diskussionsfassung)
%
% Forschungsprogramm:
% Ontophänomenalismus (OP)
% Projektive Strukturtheorie (PST)
%
% Dieses Dokument besitzt bewusst den Status eines
% Diskussions- und Forschungsdokuments.
%
% Es stellt KEIN Bestandteil der formalen Theorie dar.
%
% Ziel ist die Sammlung, Konsolidierung und Bewertung
% möglicher theoretischer Entwicklungen vor ihrer
% Überführung in DPC.tex.
%
% ==============================================================

\documentclass[11pt,a4paper]{article}

% --------------------------------------------------------------
% Pakete
% --------------------------------------------------------------

\usepackage[utf8]{inputenc}
\usepackage[T1]{fontenc}
\usepackage[german]{babel}

\usepackage{amsmath}
\usepackage{amssymb}
\usepackage{amsthm}

\usepackage{geometry}
\geometry{margin=2.5cm}

\usepackage{parskip}

\usepackage{enumitem}
\setlist{
itemsep=0.2em,
topsep=0.2em
}

\usepackage{booktabs}
\usepackage{xcolor}

\usepackage{hyperref}

\hypersetup{
colorlinks=true,
linkcolor=blue,
urlcolor=blue,
pdftitle={Diskussionspapier Projektive Chronogenese},
pdfauthor={Forschungsprogramm Ontophänomenalismus}
}

% --------------------------------------------------------------
% Theorem-Umgebungen
% --------------------------------------------------------------

\newtheorem{definition}{Definition}
\newtheorem{hypothesis}{Arbeitshypothese}
\newtheorem{remark}{Bemerkung}

% --------------------------------------------------------------
% Titel
% --------------------------------------------------------------

\title{

\textbf{Diskussionspapier}

\vspace{0.4cm}

\textbf{Zeit, Entropie und informationelle Transitionen}

\vspace{0.4cm}

Vorbereitungsdokument für das Modul

\textbf{Die Projektive Chronogenese (DPC)}

\vspace{0.4cm}

Version 2.0

}

\author{

Forschungsprogramm Ontophänomenalismus

}

\date{Juli 2026}

% ==============================================================
\begin{document}

\maketitle

% ==============================================================

\begin{abstract}

Dieses Dokument besitzt den Charakter eines Forschungs- und
Diskussionspapiers.

Es dient ausdrücklich \textbf{nicht} der Einführung neuer
Axiome oder verbindlicher Bestandteile des
Ontophänomenalismus (OP).

Sein Ziel besteht darin, mögliche theoretische Erweiterungen
für das zukünftige Modul
\textbf{Die Projektive Chronogenese (DPC)}
zusammenzuführen, zu strukturieren und kritisch zu bewerten.

Im Mittelpunkt steht die Frage, ob physikalische Zeit
nicht als fundamentale Größe verstanden werden muss,
sondern als emergente Ordnung informationeller
Transitionen innerhalb der projektiven Struktur.

Das Dokument besitzt daher bewusst einen explorativen
Charakter und formuliert zahlreiche Aussagen als
Arbeitshypothesen im Sinne der Methodologischen
Konventionen des Ontophänomenalismus (MKO).

\end{abstract}

\tableofcontents

\newpage

% ==============================================================
\section{Zweck dieses Dokuments}
% ==============================================================

Dieses Dokument verfolgt einen anderen Zweck als die
bestehenden Module des Ontophänomenalismus.

Es entwickelt keine abgeschlossene Theorie.

Es formuliert keine endgültigen mathematischen Definitionen.

Ebenso wenig beansprucht es bereits den Status eines
Bestandteils der PST.

Vielmehr dient es als gemeinsame Arbeitsgrundlage für die
spätere Entwicklung des Moduls

\begin{center}

\textbf{Die Projektive Chronogenese (DPC)}

\end{center}

Alle hier formulierten Ideen besitzen deshalb zunächst den
Status eines offenen Forschungsprogramms.

Die Aufgabe dieses Dokuments besteht insbesondere darin,

\begin{itemize}

\item bestehende Überlegungen zu sammeln,

\item unterschiedliche Hypothesen gegeneinander
abzuwägen,

\item offene Probleme sichtbar zu machen,

\item Anschlüsse an bestehende Physik zu diskutieren,

\item und eine konsolidierte Grundlage für DPC zu schaffen.

\end{itemize}

% ==============================================================
\section{Methodologischer Status}
% ==============================================================

Gemäß den Methodologischen Konventionen des
Ontophänomenalismus (MKO) wird zwischen rekonstruktiven
und postulativen Aussagen unterschieden.

Daher besitzen sämtliche zentralen Aussagen dieses
Dokuments einen expliziten Status.

\subsection*{(R) Rekonstruktiv}

Eine rekonstruktive Aussage versucht, bereits eingeführte
Bestandteile des Ontophänomenalismus logisch
weiterzuentwickeln.

Sie erhebt keinen Anspruch auf neue ontologische
Grundannahmen.

\subsection*{(P) Postulativ}

Eine postulative Aussage führt eine neue theoretische
Möglichkeit ein.

Sie besitzt deshalb zunächst ausschließlich den Charakter
einer Arbeitshypothese.

Ihre spätere Übernahme in den OP oder die PST setzt
eine eigenständige Begründung voraus.

% ==============================================================
\section{Ausgangspunkt}
% ==============================================================

Eine der ältesten Fragen der Philosophie lautet:

\begin{quote}

Was ist Zeit?

\end{quote}

Bemerkenswert ist dabei, dass nahezu alle klassischen
Antworten letztlich auf Veränderungen Bezug nehmen.

Bereits Aristoteles beschreibt Zeit als Zahl der Bewegung.

Leibniz versteht Zeit relational.

Mach kritisiert den absoluten Zeitbegriff.

Barbour spricht von einer zeitlosen Wirklichkeit.

Rovelli interpretiert Zeit als relationale Größe.

Trotz erheblicher Unterschiede besitzen diese Ansätze
eine gemeinsame Grundintuition:

\begin{quote}

Ohne Veränderung scheint Zeit ihren Sinn zu verlieren.

\end{quote}

Auch innerhalb des Ontophänomenalismus wurde bereits
früh angenommen,

dass Raum,

Zeit,

und Kausalität

nicht fundamental sind.

Sie entstehen vielmehr erst durch die projektive
Beschreibung ontischer Strukturen.

Diese Grundidee bildet den Ausgangspunkt der folgenden
Überlegungen.

% ==============================================================
\section{Die Ausgangsfrage}
% ==============================================================

Die eigentliche Fragestellung dieses Dokuments lautet nicht

\begin{quote}

Was ist Zeit?

\end{quote}

sondern vielmehr

\begin{quote}

Unter welchen Bedingungen erscheint einem
Beobachter überhaupt eine zeitliche Ordnung?

\end{quote}

Diese Umformulierung besitzt erhebliche Konsequenzen.

Denn sie verschiebt den Fokus

von der Existenz der Zeit

auf die Bedingungen ihrer Emergenz.

Damit wird Zeit nicht länger als primitives Objekt
behandelt.

Sie wird selbst erklärungsbedürftig.

Genau hierin unterscheidet sich die Projektive
Chronogenese von vielen etablierten physikalischen
Ansätzen.

% ==============================================================
\section{Abgrenzung gegenüber bestehenden Ansätzen}
% ==============================================================

Die Projektive Chronogenese beginnt nicht innerhalb der
bekannten Physik.

Sie setzt vielmehr unterhalb der bekannten Physik an.

Viele gegenwärtige Arbeiten beschäftigen sich mit
sogenannter

\begin{itemize}

\item relationaler Zeit,

\item thermodynamischer Zeit,

\item quantenmechanischer interner Zeit,

\item oder gravitativer Zeit.

\end{itemize}

Diese Ansätze besitzen eine gemeinsame Eigenschaft:

Sie setzen die mathematische Struktur der heutigen
Physik bereits voraus.

Die DPC verfolgt dagegen eine andere Fragestellung.

Sie untersucht,

wie überhaupt erstmals

eine geordnete zeitliche Beschreibung

aus einem zeitlosen ontischen Hintergrund
hervorgehen könnte.

Damit versteht sich die DPC ausdrücklich nicht als
Alternative zu bestehenden physikalischen Modellen.

Sie versteht sich vielmehr als mögliche ontologische
Grundlegung,
auf der solche Modelle später rekonstruiert werden
könnten.

% ==============================================================
% Ende Teil 1
% ==============================================================

% ==============================================================
\section{Der zeitlose Ontophänomenale Phasenraum}
% ==============================================================

Die Projektive Strukturtheorie (PST) beschreibt den
Ontophänomenalen Phasenraum (OPP) als die primäre
ontologische Struktur.

Dabei handelt es sich ausdrücklich nicht um einen
physikalischen Raum.

Ebenso wenig besitzt der OPP bereits

\begin{itemize}

\item eine Metrik,

\item eine ausgezeichnete Zeitrichtung,

\item einen physikalischen Abstand,

\item oder eine Raumzeit im Sinne der Relativitätstheorie.

\end{itemize}

Der OPP stellt vielmehr die Gesamtheit aller ontischen
Strukturen dar.

Zeit existiert auf dieser Ebene noch nicht.

Die DPC übernimmt diese Grundannahme vollständig.

\begin{remark}

Die DPC erzeugt keine neue Ontologie.

Sie untersucht ausschließlich die Bedingungen,
unter denen aus dem OPP eine Beschreibung entstehen
könnte, die wir als Zeit interpretieren.

\end{remark}

% ==============================================================
\section{Veränderung als Ausgangspunkt}
% ==============================================================

Nahezu jede Theorie der Zeit beginnt implizit mit
Veränderungen.

Auch die DPC übernimmt diesen Ausgangspunkt.

Eine Veränderung lässt sich zunächst rein formal als
Übergang

\[
S_i
\longrightarrow
S_{i+1}
\]

zwischen zwei unterscheidbaren Zuständen auffassen.

Dabei wird bewusst noch keine physikalische Bedeutung
angenommen.

Insbesondere bedeutet dies noch nicht,

dass

\[
S_i
\]

und

\[
S_{i+1}
\]

zu verschiedenen Zeitpunkten existieren.

Genau dies soll später erst erklärt werden.

% ==============================================================
\section{Von Zuständen zu Transitionen}
% ==============================================================

Der Begriff

\textbf{Zustandswechsel}

beschreibt zunächst lediglich,

dass zwei Zustände verschieden sind.

Für die DPC genügt dies jedoch nicht.

Entscheidend ist vielmehr der Vorgang,

durch den eine Struktur in eine andere übergeht.

Diesen Vorgang bezeichnet die DPC als

\begin{center}

\textbf{informationelle Transition}

\end{center}

Der Begriff

Transition

wird deshalb im weiteren Verlauf bevorzugt verwendet.

Er bezeichnet ausschließlich den elementaren
Informationsübergang zwischen zwei unterscheidbaren
Strukturen.

Er besitzt zunächst keinerlei thermodynamische,
physikalische oder psychologische Bedeutung.

% ==============================================================
\section{Arbeitshypothese: Elementare Transitionen}
% ==============================================================

\begin{hypothesis}[Elementare Transition --- Typ (P)]

Es könnte eine minimale informationelle Änderung

\[
\Delta I_{\min}
\]

existieren,

die den kleinsten projektiv unterscheidbaren Übergang
zwischen zwei Attraktorzuständen beschreibt.

\end{hypothesis}

Diese Hypothese besitzt ausdrücklich postulativen
Charakter.

Sie wird weder aus den OP-Axiomen noch aus der PST
abgeleitet.

Sie dient ausschließlich als mögliche Grundlage
für die weitere Entwicklung der DPC.

% ==============================================================
\section{Diskret oder kontinuierlich?}
% ==============================================================

An dieser Stelle entsteht unmittelbar eine grundlegende
Frage.

Sind solche Transitionen kontinuierlich

oder

diskret?

Die heutige Physik liefert hierzu kein eindeutiges Bild.

Zahlreiche Größen erscheinen gequantelt,

beispielsweise

\begin{itemize}

\item Energie,

\item Spin,

\item Drehimpuls,

\item elektrische Ladung,

\item Photonenzahl.

\end{itemize}

Demgegenüber werden Raum und Zeit innerhalb der
Relativitätstheorie als kontinuierliche Mannigfaltigkeit
modelliert.

Die DPC muss deshalb beide Möglichkeiten zunächst
offenhalten.

% ==============================================================
\section{Die Planck-Zeit}
% ==============================================================

Oft wird argumentiert,

die Existenz der Planck-Zeit

\[
t_P
=
\sqrt{\frac{\hbar G}{c^5}}
\]

beweise bereits,

dass Zeit diskret sei.

Diese Schlussfolgerung erscheint jedoch nicht zwingend.

Die Planck-Zeit beschreibt zunächst lediglich eine
natürliche Kombination fundamentaler Konstanten.

Ob sie tatsächlich einer minimalen physikalischen
Zeitauflösung entspricht,

ist bislang experimentell ungeklärt.

Die DPC verwendet die Planck-Zeit deshalb nicht als
Argument für eine fundamentale Diskretheit.

% ==============================================================
\section{Variante A und Variante B}
% ==============================================================

Im Verlauf der Diskussion entstanden zwei prinzipielle
Möglichkeiten.

\subsection*{Variante A}

Der OPP selbst wäre bereits diskret.

Diskrete Zeit ergäbe sich unmittelbar aus der Struktur
des ontischen Phasenraums.

\subsection*{Variante B}

Der OPP bleibt vollständig kontinuierlich.

Die Diskretheit entsteht ausschließlich durch den
Projektionsprozess.

Beide Möglichkeiten wurden ausführlich diskutiert.

% ==============================================================
\section{Entscheidung zugunsten von Variante B}
% ==============================================================

\begin{hypothesis}[Kontinuierlicher OPP --- Typ (P)]

Für die weitere Entwicklung der DPC wird zunächst
Variante B als bevorzugte Arbeitshypothese verwendet.

Der Ontophänomenale Phasenraum bleibt kontinuierlich.

Diskrete Zeit entsteht ausschließlich als Folge der
projektiven Beschreibung.

\end{hypothesis}

Diese Entscheidung besitzt mehrere Vorteile.

Sie verhindert insbesondere,

dass bereits auf ontischer Ebene

eine ausgezeichnete Gitterstruktur eingeführt werden
müsste.

Damit bleiben die Symmetrien des OPP möglichst allgemein.

% ==============================================================
\section{Lorentz-Invarianz}
% ==============================================================

Die Entscheidung für einen kontinuierlichen OPP erfolgt
nicht,

weil dadurch automatisch die Lorentz-Invarianz bewiesen
würde.

Ein solcher Beweis existiert an dieser Stelle nicht.

Vielmehr handelt es sich um eine methodologische
Überlegung.

Ein fundamentaler Gitterraum würde sehr leicht ein
bevorzugtes Bezugssystem einführen.

Ein kontinuierlicher OPP vermeidet diese Schwierigkeit
zunächst.

Die eigentliche Rekonstruktion der Lorentz-Symmetrie
bleibt jedoch Aufgabe der PST.

\begin{remark}

Die DPC behauptet daher ausdrücklich nicht,

dass Variante B die Lorentz-Invarianz bereits erklärt.

Sie schafft lediglich günstigere Voraussetzungen für
eine spätere Rekonstruktion.

\end{remark}

% ==============================================================
\section{Projektive Auflösung}
% ==============================================================

Falls Variante B zutrifft,

besitzt nicht der OPP selbst eine minimale Struktur,

sondern der Projektionsprozess.

Die kleinste unterscheidbare Änderung wäre dann keine
Eigenschaft der Realität,

sondern eine Eigenschaft der Beschreibung.

Damit erhält

\[
\Delta I_{\min}
\]

eine neue Interpretation.

Es beschreibt nicht die kleinste ontische Veränderung,

sondern die kleinste durch den Projektionsoperator

\[
\mathcal P^\ast
\]

noch unterscheidbare Änderung.

Dies verschiebt die gesamte Diskussion

von einer ontologischen

zu einer projektiven Fragestellung.

% ==============================================================
% Ende Teil 2
% ==============================================================

% ==============================================================
\section{Zeit als Zählgröße informationeller Transitionen}
% ==============================================================

Unter der Annahme einer minimal projektiv unterscheidbaren
Transition könnte Zeit vollständig neu interpretiert werden.

Zeit wäre dann keine fundamentale Koordinate,
sondern eine Zählgröße.

Es gelte

\[
N_{\mathrm{trans}}
=
\text{Anzahl projektiv unterscheidbarer Transitionen}.
\]

Dann könnte die physikalische Zeit als

\[
t
=
\tau \cdot N_{\mathrm{trans}}
\]

geschrieben werden.

Dabei besitzt

\[
\tau
\]

keinen ontischen Status.

Es beschreibt lediglich die physikalische Interpretation
einer einzelnen projektiven Transition.

Damit wäre Zeit keine unabhängige Größe,
sondern ausschließlich eine Ordnung der Transitionen.

% ==============================================================
\section{Informationeller Zeitpfeil}
% ==============================================================

Die DPC unterscheidet streng zwischen

\begin{enumerate}

\item einer bloßen Reihenfolge von Transitionen,

\item und einem gerichteten Zeitpfeil.

\end{enumerate}

Eine Reihenfolge allein erzeugt noch keine
Irreversibilität.

Der Zeitpfeil verlangt zusätzliche Bedingungen.

Insbesondere muss erklärt werden,

weshalb makroskopische Prozesse eine bevorzugte Richtung
besitzen,

obwohl die zugrunde liegenden mikroskopischen
Beschreibungen häufig zeitumkehrsymmetrisch erscheinen.

Die DPC betrachtet den Zeitpfeil deshalb nicht
als Fundament,

sondern als emergente Eigenschaft großer Mengen
informationeller Transitionen.

% ==============================================================
\section{Entropie als emergente Statistik}
% ==============================================================

In der klassischen Thermodynamik gilt

\[
S
=
k_B \ln W.
\]

Hier beschreibt

\[
W
\]

die Anzahl möglicher Mikrozustände.

Die DPC schlägt eine andere ontologische Einordnung vor.

\begin{hypothesis}[Emergente Entropie --- Typ (P)]

Entropie könnte keine fundamentale Eigenschaft des OPP
sein,

sondern erst auf der projektiven Ebene
als makroskopische Statistik vieler Transitionen
entstehen.

\end{hypothesis}

Damit verschiebt sich die Kausalrichtung.

Nicht die Entropie erzeugt den Zeitpfeil,

sondern die projektive Dynamik erzeugt zunächst
Transitionen,

deren statistische Beschreibung später
als Entropie erscheint.

% ==============================================================
\section{Die Kausalkette der DPC}
% ==============================================================

Unter dieser Interpretation ergibt sich folgende
Arbeitshypothese:

\[
\boxed{
\Delta I_{\min}
\Longrightarrow
\text{Transition}
\Longrightarrow
\text{Zeit}
\Longrightarrow
\text{Entropie}
}
\]

Diese Reihenfolge besitzt ausschließlich
heuristischen Charakter.

Sie stellt keine bewiesene Ableitung dar,

sondern beschreibt die gegenwärtige
Forschungsrichtung der DPC.

% ==============================================================
\section{Zusammenhang mit der Ontologischen Attraktoren-Dynamik}
% ==============================================================

Die Ontologische Attraktoren-Dynamik (OAD)
führt stabile Attraktoren als fundamentale
Strukturelemente der PST ein.

Die DPC interpretiert diese Attraktoren
nun dynamisch.

Ein stabiler Attraktor bedeutet gerade nicht,

dass keinerlei Veränderung stattfindet.

Vielmehr könnte ein Attraktor permanent
mikroskopische Transitionen durchlaufen,

ohne seine makroskopische Identität zu verlieren.

Stabilität würde somit bedeuten,

dass sich die projektiv beobachtbare Struktur
trotz kontinuierlicher interner Änderungen erhält.

% ==============================================================
\section{Heuristisches Stabilitätsmaß}
% ==============================================================

Als erste Arbeitshypothese wurde vorgeschlagen,

die Stabilität eines Attraktors durch

\[
\Sigma_A
=
\frac
{N_{\mathrm{stabil}}}
{N_{\mathrm{gesamt}}}
\]

zu beschreiben.

Dabei bezeichnet

\[
N_{\mathrm{stabil}}
\]

die Anzahl erfolgreich überstandener Transitionen,

während

\[
N_{\mathrm{gesamt}}
\]

alle beobachteten Transitionen umfasst.

\begin{definition}[Heuristisches Modell]

Das Stabilitätsmaß

\[
\Sigma_A
\]

ist gegenwärtig ausschließlich
ein Modell.

Es besitzt noch keinen ontologischen Status
und kann im Verlauf der PST
durch geeignetere Größen ersetzt werden.

\end{definition}

% ==============================================================
\section{Numerische Perspektive}
% ==============================================================

Die bisherigen PST-Simulationen verwendeten bereits

\begin{itemize}

\item Diffusion Maps,

\item Kernel PCA,

\item Eigenwertspektren,

\item Gram-Matrizen.

\end{itemize}

Hieraus ergibt sich unmittelbar
eine neue Forschungsfrage.

Kann die Stabilität eines Attraktors

direkt

aus dem Eigenwertspektrum rekonstruiert werden?

Insbesondere erscheint folgende Vermutung
interessant:

\[
\Sigma_A
=
F(\lambda_i).
\]

Dabei bezeichnen

\[
\lambda_i
\]

die dominanten Eigenwerte
der Diffusionsoperatoren.

Eine solche Beziehung wäre numerisch überprüfbar
und könnte den ersten direkten Test
der DPC ermöglichen.

% ==============================================================
\section{Beziehung zum Fragedifferential}
% ==============================================================

Im Modul

DFD.tex

wurde das Fragedifferential

\[
\mathcal D(E)
\]

als Maß informationeller Differenz
zwischen Erklärungsräumen eingeführt.

Die DPC führt nun

\[
\Delta I_{\min}
\]

ein.

Beide Größen besitzen ähnliche Rollen,
sind jedoch nicht identisch.

\begin{remark}

Nach gegenwärtigem Stand erscheint
es plausibel,

dass

\[
\Delta I_{\min}
\]

als lokale projektive Minimalform
des Fragedifferentials interpretiert
werden könnte.

Ob diese Vermutung mathematisch tragfähig ist,

bleibt Gegenstand zukünftiger Untersuchungen.

\end{remark}

% ==============================================================
\section{Mögliche mathematische Definition von
\texorpdfstring{$\Delta I_{\min}$}{Delta I}}
% ==============================================================

Gegenwärtig bleibt

\[
\Delta I_{\min}
\]

bewusst offen.

Mehrere Kandidaten erscheinen möglich.

Beispielsweise

\begin{itemize}

\item KL-Divergenz,

\item Jensen-Shannon-Divergenz,

\item Mutual Information,

\item Wasserstein-Distanzen,

\item projektive Eigenwertabstände.

\end{itemize}

Die DPC legt sich bewusst noch
auf keine dieser Möglichkeiten fest.

Die endgültige Wahl soll sich
an mathematischer Konsistenz
und numerischer Brauchbarkeit orientieren.

% ==============================================================
\section{Beziehungen zu bestehenden Forschungsprogrammen}
% ==============================================================

Die DPC versteht sich ausdrücklich
nicht als Konkurrenz,

sondern als möglicher metatheoretischer Rahmen.

Insbesondere bestehen interessante Berührungspunkte zu

\begin{itemize}

\item Julian Barbour,

\item Carlo Rovelli,

\item John Wheeler ("It from Bit"),

\item Landauer,

\item Verlinde,

\item Informationsphysik,

\item Relational Quantum Mechanics.

\end{itemize}

Ebenso bestehen Überschneidungen
mit verschiedenen Ansätzen
diskreter Raumzeitmodelle.

Der entscheidende Unterschied besteht jedoch darin,

dass die DPC

die Diskretheit nicht notwendigerweise
dem OPP zuschreibt,

sondern dem Projektionsprozess.

% ==============================================================
\section{Aktuelle Forschungsperspektiven}
% ==============================================================

Neuere theoretische Arbeiten zur Entstehung
von Zeit,

Informationsfluss,

Quantenkausalität
und thermodynamischer Emergenz
zeigen,

dass die Frage nach einer
informationellen Rekonstruktion
der Zeit zunehmend an Bedeutung gewinnt.

Die DPC versteht sich deshalb
nicht als isolierter Ansatz,

sondern als möglicher Beitrag
zu einem größeren Forschungsfeld
über die Emergenz physikalischer Zeit.

% ==============================================================
\section{Vorläufiges Gesamtfazit}
% ==============================================================

Die hier entwickelten Überlegungen
bilden ausdrücklich

kein fertiges Theoriegebäude.

Sie definieren vielmehr
das Forschungsprogramm,
auf dessen Grundlage
das spätere Modul

\texttt{DPC.tex}

entwickelt werden soll.

Die zentrale Leitidee lautet:

\begin{quote}

Zeit könnte weder fundamental
noch ontologisch primitiv sein.

Sie könnte vielmehr vollständig
als Folge projektiv unterscheidbarer
informationeller Transitionen
rekonstruiert werden.

\end{quote}

Sollte sich diese Arbeitshypothese
als tragfähig erweisen,

würde die DPC
das zeitliche Gegenstück
zur Projektiven Reduktion (DPR)
bilden

und damit einen weiteren Baustein
für die vollständige Rekonstruktion
von Raum, Zeit und Physik
innerhalb der Projektiven Strukturtheorie liefern.

% ==============================================================
% Ende Teil 3
% ==============================================================

% ==============================================================
\appendix
% ==============================================================

\section{Offene Forschungsfragen}
\label{app:fragen}

Die DPC befindet sich gegenwärtig im Status eines
Forschungsprogramms. Zahlreiche Fragen sind bewusst
offengehalten und sollen erst im Rahmen der weiteren
Entwicklung der PST beantwortet werden.

\subsection*{(F1) Was ist eine informationelle Transition?}

Bislang wurde der Begriff ausschließlich heuristisch
eingeführt.

Es bleibt zu klären,

\begin{itemize}

\item ob Transitionen mathematisch primitive Objekte sind,

\item oder aus tieferliegenden Strukturen rekonstruiert
werden können.

\end{itemize}

% --------------------------------------------------------------

\subsection*{(F2) Wie wird \texorpdfstring{$\Delta I_{\min}$}{Delta I} formal definiert?}

Mehrere Kandidaten erscheinen derzeit plausibel.

Unter anderem

\begin{itemize}

\item KL-Divergenz,

\item Jensen-Shannon-Divergenz,

\item Mutual Information,

\item Wasserstein-Distanz,

\item projektive Eigenwertabstände,

\item andere geometrische Informationsmaße.

\end{itemize}

Die Auswahl soll ausschließlich
auf mathematischer Konsistenz,
empirischer Anschlussfähigkeit
und numerischer Stabilität beruhen.

% --------------------------------------------------------------

\subsection*{(F3) Existiert eine minimale Transition überhaupt?}

Möglicherweise existiert überhaupt keine
fundamentale Minimalgröße.

Ebenso denkbar ist,

dass

\[
\Delta I_{\min}
\]

ausschließlich durch die endliche
Auflösung des Projektionsoperators entsteht.

Auch diese Möglichkeit bleibt offen.

% --------------------------------------------------------------

\subsection*{(F4) Wie entsteht der Zeitpfeil?}

Die DPC beschreibt zunächst lediglich
eine geordnete Folge von Transitionen.

Eine Richtung der Zeit
ergibt sich daraus noch nicht.

Zu untersuchen bleibt insbesondere,

\begin{itemize}

\item welche Rolle statistische Irreversibilität spielt,

\item ob der Zeitpfeil ausschließlich projektiv entsteht,

\item oder ob zusätzliche Selektionsprinzipien erforderlich
sind.

\end{itemize}

% --------------------------------------------------------------

\subsection*{(F5) Zusammenhang mit der Thermodynamik}

Kann der zweite Hauptsatz

vollständig

als emergente Folge projektiver
Informationsverluste verstanden werden?

Oder bleiben zusätzliche physikalische
Annahmen notwendig?

Diese Frage besitzt hohe Priorität.

% --------------------------------------------------------------

\subsection*{(F6) Numerische Testbarkeit}

Die DPC gewinnt erheblich an Stärke,

wenn ihre Aussagen numerisch überprüfbar werden.

Insbesondere sollen untersucht werden:

\begin{itemize}

\item Stabilitätsmaße auf Attraktoren,

\item Eigenwertspektren,

\item Diffusion Maps,

\item Kernelmethoden,

\item projektive Informationsmetriken.

\end{itemize}

% ==============================================================
\section{Methodologische Einordnung}
% ==============================================================

Gemäß den Methodologischen Konventionen (MKO)

besitzt dieses Dokument ausschließlich
Diskussionscharakter.

Es führt

\begin{itemize}

\item keine neuen OP-Axiome,

\item keine bindenden mathematischen Definitionen,

\item keine endgültigen ontologischen Festlegungen

\end{itemize}

ein.

Sämtliche neu eingeführten Begriffe sind

bis zur späteren Ausarbeitung von

\texttt{DPC.tex}

als

heuristische Forschungsobjekte

zu verstehen.

% ==============================================================
\section{Einordnung innerhalb der Dokumentstruktur}
% ==============================================================

Die geplante Position der Projektiven Chronogenese
innerhalb des OP/PST ergibt sich derzeit wie folgt:

\bigskip

\begin{center}

OP

$\downarrow$

OPP

$\downarrow$

OAD

$\downarrow$

DPR

$\downarrow$

\boxed{\textbf{DPC}}

$\downarrow$

DPD

$\downarrow$

PST

\end{center}

\bigskip

Dabei übernimmt

\begin{itemize}

\item DPR die projektive Emergenz von Raum,

\item DPC die projektive Emergenz der Zeit,

\item DPD die projektive Emergenz der
Subjekt-Objekt-Differenz.

\end{itemize}

In dieser Struktur bildet DPC
das zeitliche Gegenstück
zur Projektiven Reduktion.

% ==============================================================
\section{Abschließende Bemerkung}
% ==============================================================

Die in diesem Dokument entwickelten Ideen entstanden
im Rahmen eines mehrmonatigen iterativen
Forschungsprozesses.

Sie wurden fortlaufend

\begin{itemize}

\item logisch überprüft,

\item mathematisch hinterfragt,

\item philosophisch reflektiert,

\item mit numerischen PST-Überlegungen abgeglichen,

\item sowie durch den strukturierten Dialog
zwischen menschlichem Autor
und mehreren unabhängigen KI-Systemen
kritisch diskutiert.

\end{itemize}

Die vorliegende Fassung versteht sich daher
nicht als Abschluss,

sondern als gesicherter Zwischenstand.

Sie soll den vollständigen konzeptionellen Kontext
dauerhaft dokumentieren,

damit die spätere Ausarbeitung von

\texttt{DPC.tex}

auf einer konsolidierten Grundlage erfolgen kann,
ohne dass wesentliche Überlegungen,
Diskussionen oder Argumentationslinien
verlorengehen.

\bigskip

\begin{center}
\Large
\textit{Ende des Diskussionspapiers}
\end{center}

\end{document}

